\documentclass[aspectratio=169]{beamer}

% --- Tema y colores UTEQ ---
\usetheme{Madrid}
\definecolor{uteqAzul}{RGB}{0,48,135}
\definecolor{uteqNaranja}{RGB}{245,166,35}
\setbeamercolor{structure}{fg=uteqAzul}
\setbeamercolor{palette primary}{bg=uteqAzul,fg=white}
\setbeamercolor{palette secondary}{bg=uteqNaranja,fg=white}
\setbeamercolor{frametitle}{bg=uteqAzul,fg=white}
\setbeamercolor{title}{fg=uteqAzul}

% --- Paquetes ---
\usepackage[utf8]{inputenc}
\usepackage[spanish]{babel}
\usepackage{amsmath}
\usepackage{hyperref}

% --- Logo UTEQ (descomentar cuando tengas el archivo logo_uteq.png en este directorio) ---
% \usepackage{graphicx}
% \logo{\includegraphics[height=0.5cm]{logo_uteq.png}}

% --- Metadata (Selección de materiales con método Ashby y Profesor son sustituidos automáticamente por Mitzlia) ---
\title{Selección de materiales con método Ashby}
\author{Profesor}
\institute{UTEQ}
\date{\today}

\begin{document}

\begin{frame}
  \titlepage
\end{frame}

\begin{frame}{\frametitle{Portada}}
  \begin{itemize}
    \item \textbf{Título:} Selección de materiales
    \item \textbf{Subtítulo:} Aplicando el método de Ashby
    \item \textbf{Autor:} Nombre del Profesor
    \item \textbf{Institución:} UTEQ — Ingeniería Mecánica
    \item \textbf{Fecha:} 2023-11
  \end{itemize}
\end{frame}

\begin{frame}{\frametitle{Agenda}}
  \begin{itemize}
    \item Introducción a la selección de materiales
    \item Fundamentos del método de Ashby
    \item Aplicación práctica del método
    \item Caso de estudio
    \item Resumen
    \item Preguntas
  \end{itemize}
\end{frame}

\begin{frame}{\frametitle{Introducción}}
  \medskip\textbf{Conceptos clave:}
  \begin{itemize}
    \item Necesidad de seleccionar materiales óptimos
    \item Factores que influyen en la elección
  \end{itemize}
  \medskip\textbf{Ejemplo:}
  \begin{itemize}
    \item Comparación entre acero y aluminio en estructuras
  \end{itemize}
\end{frame}

\begin{frame}{\frametitle{Método de Ashby}}
  \medskip\textbf{Conceptos clave:}
  \begin{itemize}
    \item Proceso sistemático para comparar materiales
    \item Uso de propiedades normalizadas
  \end{itemize}
  \medskip\textbf{Ecuación:}
  \begin{itemize}
    \item \$\frac{\sigma}{\rho} = k\$ (relación entre resistencia y densidad)
  \end{itemize}
\end{frame}

\begin{frame}{\frametitle{Pasos del método}}
\end{frame}

\begin{frame}{\frametitle{Caso de estudio}}
  \begin{itemize}
    \item Selección de material para un componente térmico
    \item Análisis de propiedades térmicas y mecánicas
    \item Resultados del método de Ashby
  \end{itemize}
\end{frame}

\begin{frame}{\frametitle{Resumen}}
  \begin{itemize}
    \item Selección basada en criterios técnicos
    \item Importancia del método de Ashby en ingeniería
    \item Ventajas de la metodología sistemática
  \end{itemize}
\end{frame}

\begin{frame}{\frametitle{Preguntas y Discusión}}
  \begin{itemize}
    \item Espacio para consultas
    \item Reflexión final: ¿Cómo optimizar la selección de materiales?
  \end{itemize}
\end{frame}


\end{document}
